\scanpage{3}\mspage{2}
\margin{b) se déduit de a) et de 1.2. a)}
\diagram{\begin{tikzcd}[ampersand replacement=\&,column sep=small,row sep=small,
execute at end picture={\draw (aa.north west) .. controls +(0,-.25) and +(.04,.16) .. (aa.south west) .. controls +(-.04,-.10) and +(-.16,.03) .. (bb.north west) .. controls +(.08,-.10) and +(-.02,.25) .. (bb.south west);}]
{\text{(1.3)}} \& |[alias=aa]| {\text{a)}} \& {\mathcal F\hookrightarrow\mathcal P(M)} \& {F\mapsto\widetilde F=\operatorname{\text{Dépl}}(F)\subset M} \\
\& |[alias=bb]| {\text{b)}} \& {\mathcal F\hookrightarrow\mathcal P(\mathcal P(M))} \& {\text{plus précisément }\mathcal F\to\operatorname{Fig}(M)} \\
\& {\substack{\text{ne sert}\\\text{à rien}}}\arrow[u,bend left=24]\arrow[rr,bend right=13] \& \& {F\mapsto\overset{\deleted{\unsure{\sim}}}{\mathrm F}\mathrm{ig}_{M}(F)\subset\mathcal P(M)}
\end{tikzcd}}
où $\operatorname{Fig}(M)$ désigne la partie de $\mathcal P(\mathcal P(\unsure{\cdot}M))$ formée des figures ensemblistes dans $M$. \deleted{On}

On peut considérer que la première application correspond à une relation entre $M$ et $\mathcal F$, appelée relation d'incidence. Pour une figure $F$, $\widetilde F$ s'appelle l'ensemble des \uline{multistrates} \uline{incidentes}, ou le \uline{déployement} de la figure $F$. Si $X\in M$, $F\in\mathcal F$, on dira que la multistrate $X$ est \uline{incidente} à la figure $F$, ou encore que c'est une \uline{strate de} \deleted{$F$}\inserted{\uline{la figure} $F$}, si $X\in\widetilde F$. D'autre part, tout\deleted{e} élément \inserted{$X$} de $M$ (ie. \deleted{toute} multistrate), regardé comme une figure par (1.2), admet un déployement $\widetilde X$, et \deleted{la multistrate} \deleted{\unsure{mais}} \deleted{\unsure{on|ou}} on pose
\margin{$\unsure{\bullet|\cdot}$}
\begingroup\makeatletter\tagsleft@true\let\veqno\@@leqno\makeatother
\[
\begin{aligned}
\text{\unsure{``}}\operatorname{Fig}_{M}(F)&=\{\widetilde X\mid X\in\widetilde F=\operatorname{Depl}(F)\}\unsure{\subset|\in}\mathcal P(M)\\
\text{ie.}\quad \operatorname{Fig}_{M}(F)&\in\mathcal P(\mathcal P(M))
\end{aligned}\tag*{(1.4)}
\]
\endgroup
et il résultera des axiomes que c'est une figure ensembliste dans $M$, indexée fidèlement par l'\unsure{ens|un} $\widetilde F$ des strates de $F$.

\deleted{\unsure{En fait}} \deleted{$\unsure{\mathcal F|F}$} En fait, $M$ sera muni d'une relation d'ordre $\leq$, cf plus bas, et $\widetilde F\inserted{\subset M}$ sera une partie fermée de \deleted{\unsure{la|de}} $M$, et pour tout $X\in\widetilde F$, \deleted{\unsure{$X$}} on aura
\begingroup\makeatletter\tagsleft@true\let\veqno\@@leqno\makeatother
\[
\widetilde X=\{Y\in M\mid Y\leq\unsure{x|X}\}\tag*{(1.5.)}
\]
\endgroup
